\documentclass[11pt]{article}
\usepackage[a4paper,margin=25mm]{geometry}
\usepackage{amsmath,amssymb,amsthm}
\usepackage[hidelinks]{hyperref}
\newtheorem{theorem}{Theorem}
\setlength{\emergencystretch}{1em}
\begin{document}
\title{A first-order splitting with commutator coefficient one quarter}
\author{Chengzhe Gao}\date{}\maketitle
Conjecture 6551 states that the leading commutator coefficient of a
first-order fractional-step splitting is always one half. Consider
\[
 \Phi_h=\exp(3hA/4)\exp(hB)\exp(hA/4).
\]
Both component time sums are one. This is a three-stage splitting, with
positive coefficients, of the evolution generated by $A+B$.

We use the standard formal-order convention: compare the Taylor series
of the product with the exact exponential in noncommuting operators;
see \cite{BCM}, Section 3.

\begin{theorem}
The method has formal order exactly one and
\[
 \Phi_h-\exp(h(A+B))=\frac{h^2}{4}[A,B]+O(h^3).
\]
In particular the claimed coefficient one half is not universal.
\end{theorem}
\begin{proof}
Multiplying the three exponential expansions through degree two gives
\[
 \Phi_h=I+h(A+B)+h^2\left(
 \frac{A^2}{2}+\frac{3AB}{4}+\frac{BA}{4}+\frac{B^2}{2}\right)+O(h^3).
\]
The corresponding coefficient in $\exp(h(A+B))$ is
$(A^2+AB+BA+B^2)/2$. Their difference is $[A,B]/4$.
The constant and linear terms agree, while the coefficient of the word
$AB$ differs, so the universal formal order is exactly one, not two.
For a concrete noncommuting example take
\[
 A=\begin{pmatrix}0&1\\0&0\end{pmatrix},\qquad
 B=\begin{pmatrix}0&0\\1&0\end{pmatrix}.
\]
Then $[A,B]=\operatorname{diag}(1,-1)\ne0$, and the second-order local
error is nonzero. These finite matrices also give the usual analytic
Taylor interpretation. Reversing the error-sign convention changes the
coefficient to $-1/4$, whose magnitude is still not one half.
\end{proof}

The source defines first-order splitting as a first-order fractional-step
approximation; it does not restrict it to the two-factor Lie--Trotter
formula. The three-stage method above has exactly that order. We do not
count a higher-order method as first-order, and we do not claim a different
coefficient for the particular two-factor Lie--Trotter formula.

\newpage
\paragraph{Formal verification.}
Lean represents a noncommutative word by a list of two letters and a formal
series by its rational coefficient at every word. The exponential of
$ahA$ has coefficient $a^n/n!$ on $A^n$ and zero on all other words;
the analogous definition is used for $B$. Multiplication sums over all
splits of a word into a prefix and suffix. The file proves that these
splits reconstruct the word and that every concatenation occurs. Thus
the object checked is the actual exponential product as an infinite
formal series, not an unrelated table of error constants.

The finite checking list contains every word of length at most four.
A structural proof on arbitrary lists proves its completeness. The
checking lemma connects finite certificates to equality at every word
below the degree bound. All words of degrees zero and one match the exact
series, every word through degree two has the asserted commutator error,
and the word $AB$ proves failure of order two. Theorem
\texttt{conjecture6551\_false} negates the claimed universal coefficient
for methods of exactly order one.

The additional substitution lemma proves that coefficient equality
preserves every homogeneous finite evaluation in arbitrary target
algebras. This covers substitution of arbitrary matrices or operators;
the argument never assumes $AB=BA$. The concrete two-by-two matrices,
their noncommutation, and every entry of the degree-two error obtained
from all four words are also checked. The project uses Lean 4.19.0,
\texttt{Std}, and its exact rational arithmetic. No admitted proofs,
new axioms, or external decision oracle are used.
\begin{thebibliography}{1}
\bibitem{BCM} S. Blanes, F. Casas and A. Murua,
\emph{Splitting and composition methods in the numerical integration
of differential equations}, 2008.
\href{https://arxiv.org/abs/0812.0377}{arXiv:0812.0377}.
\end{thebibliography}
\end{document}
